\chapter*{Introduction}
\addcontentsline{toc}{chapter}{Introduction}

\section*{Why these notes?}

Superconducting circuits have, over the last two decades, become one of
the leading hardware platforms for quantum computation, quantum simulation,
and the experimental study of quantum optics in the microwave regime.
Working in this field requires fluency in several different bodies of
knowledge: the postulates and basic toolbox of non-relativistic quantum
mechanics, the methods of quantum optics (originally developed for
atom--photon interactions), and circuit theory adapted to the quantum
regime.

Many excellent textbooks cover each of these subjects individually.
The purpose of these notes is more modest: to collect, in one place
and in a uniform style, the specific subset of results that one
typically uses on a daily basis when designing, simulating or analysing
a superconducting-qubit experiment. The text is therefore intentionally
not exhaustive --- formal proofs, edge cases and historical context are
mostly omitted in favour of clarity and density. Where a derivation
is given, it is given in full; where it is not, a pointer to the
relevant literature is provided.

\section*{Scope and structure}

The document is organised in five parts.
\begin{itemize}
  \item \textbf{Part~I --- Quantum mechanics foundations.} Postulates,
        Hilbert space and operators, the harmonic oscillator and
        two-level system, density matrices and composite systems,
        time-dependent perturbation theory, and the rotating-wave
        approximation.
  \item \textbf{Part~II --- Open quantum systems.} The Lindblad master
        equation, $T_1$ and $T_2$, input--output theory, and the
        quantum Zeno effect.
  \item \textbf{Part~III --- Quantum optics.} Quantisation of the
        electromagnetic field, Fock, coherent and squeezed states, the
        Jaynes--Cummings model, dispersive coupling, the AC Stark
        shift, and the Purcell effect.
  \item \textbf{Part~IV --- Superconducting circuits.} Circuit
        quantisation, the Josephson effects, the transmon and related
        qubits, parametric and four-wave-mixing processes, dispersive
        readout, single- and two-qubit gates, and the standard
        calibration experiments (resonator and qubit spectroscopy,
        Rabi, Ramsey).
  \item \textbf{Part~V --- Quantum error correction.} Codes, the
        stabiliser formalism, the surface code and a brief look at
        bosonic codes.
\end{itemize}

\section*{Pedagogical convention}

Every topic follows a four-step structure that we adopt deliberately
throughout the document:
\begin{enumerate}
  \item \emph{Phenomenon.} A short description, in words, of what is
        observed and why it deserves attention.
  \item \emph{Statement.} The result, stated formally as a theorem,
        proposition, definition or postulate.
  \item \emph{Derivation.} The mathematical argument that establishes
        the statement.
  \item \emph{Example / application.} A pointer (sometimes worked out,
        sometimes not) to a context in which the result is used.
\end{enumerate}
A reader who already knows the topic can scan the statements and skip
the rest; a reader meeting the topic for the first time can build
intuition from the phenomenon, then follow the derivation, then see
where it lands.

\section*{Notation and conventions}

We carry all factors of $\hbar$ explicitly throughout: angular
frequencies $\omega$ enter Hamiltonians as $\hbar\omega$, the
canonical commutator reads $[\hat{x},\hat{p}]=i\hbar$, time evolution
is generated by $e^{-i\hat{H}t/\hbar}$, and so on. Nothing is absorbed
into ``natural units''.

State vectors are written as $\ket{\psi}$ and their duals as
$\bra{\psi}$, with $\braket{\phi}{\psi}\in\C$ the inner product.
\emph{All operators are decorated with a hat} --- $\hat{O}$, $\hat{H}$,
$\hat{a}$, $\hat{\sigma}_x$ and so on --- whether or not a classical
counterpart could in principle be confused with them. Pauli operators
are denoted $\hat\sigma_x,\hat\sigma_y,\hat\sigma_z$ with the
convention $\hat\sigma_z\ket{0}=+\ket{0}$ and
$\hat\sigma_z\ket{1}=-\ket{1}$, so that $\ket{0}$ is the qubit ground
state. The bosonic ladder operators $\hat a,\hat a^\dagger$ satisfy
$[\hat a,\hat a^\dagger]=\id$. Throughout, $\hilbert$ denotes a
complex separable Hilbert space and $\id$ the identity operator on
it; we render the identity by the standard mathematical unit symbol
$\mathbf{1}$.

\section*{How to read this document}

Readers familiar with non-relativistic quantum mechanics may skim
Part~I, treating it as a reference for notation, and proceed directly
to Part~III. Readers new to the field are encouraged to read in order:
the parts are deliberately layered, with later parts using results,
notation, and intuition built up earlier.
